\subtituloA{Los cuatro campos del diseño web}

El diseño de un sitio web se articula en torno a cuatro campos fundamentales que determinan su éxito y funcionalidad:

\definicion{Campo visual}{Aspecto estético y de comunicación de la página web. Su objetivo principal es atraer a los clientes y causar una buena sensación en los visitantes, aumentando así las probabilidades de que vuelvan a visitar el sitio. Engloba la usabilidad, la accesibilidad y la experiencia de usuario.}

\definicion{Campo tecnológico}{Decisión sobre cómo se implementará el sitio y qué tecnologías se utilizarán. De esta elección dependen las características fundamentales del sitio: si será una página estática o dinámica, qué sistemas de bases de datos se emplearán, qué lenguajes de servidor se utilizarán y qué frameworks se aplicarán.}

\definicion{Campo de distribución}{Aspecto relativo a la velocidad y fiabilidad de la transmisión de datos entre el servidor y los clientes. Incluye la optimización del ancho de banda, el uso de redes de distribución de contenido y la garantía de disponibilidad del servicio durante las 24 horas del día.}

\definicion{Campo de propósito}{Objetivo fundamental que persigue el sitio web. En la mayoría de los casos la finalidad es económica, buscando alcanzar una gran difusión del producto en el mercado laboral, aunque también puede ser informativa, como la difusión de actos o eventos de carácter general que puedan interesar al público.}
